\documentclass[11pt,a4paper]{article}
\usepackage[margin=25mm]{geometry}
\usepackage[T1]{fontenc}
\usepackage{lmodern,amsmath,amssymb,amsthm,microtype}
\usepackage[hidelinks]{hyperref}
\hypersetup{pdftitle={Efficient submodule operations in pruning}}
\newtheorem{proposition}{Proposition}
\newcommand{\im}{\operatorname{im}}
\newcommand{\F}{\mathbb F_2}
\newcommand{\Id}{\operatorname{Id}}
\begin{document}

\section{Efficient submodule operations in pruning}
Let $M$ be a finitely presented bigraded persistence module over $\F$, let
$N=M(\delta)$ for $\delta=(2\epsilon,2\epsilon)$ with $\epsilon\ge0$, and let $\eta:M\to N$
be the canonical structure map. For the map families selected for pruning,
an I pass starts at $U\subseteq M$ and computes
\[
 U_0=U,\qquad U_i=U_{i-1}\cap f_i^{-1}(J),\qquad J=\eta U.
\]
A K pass starts at $V\subseteq M$ and computes
\[
 L=\eta V+\sum_{j=1}^{r}h_j(V),\qquad V_{\mathrm{next}}=\eta^{-1}(L).
\]
The target $J$ and source $V$ remain fixed throughout their respective
passes. The following changes reduce the cost of these operations while
preserving the resulting graded submodules.

\subsection{Projected kernel computation}
Let $F$ be a free-generator lift of $f:M\to N$, let $S$ lift generators
of $J\subseteq N$, and let $R_N$ be the target relation matrix, all in the
same free target basis. At each grade, a free-source vector $x$ represents
an element of $f^{-1}(J)$ precisely when
\[
 Fx=Sy+R_Nz
 \quad\text{for some }y,z.
\]
Consequently, if $F$ has $n$ columns, the projection onto the first $n$
coordinates gives the preimage in the free source:
\begin{equation}\label{eq:projection}
 D=[\,F\mid S\mid R_N\,],\qquad
 \pi=[\,\Id_n\mid0\mid0\,],\qquad \pi(\ker D).
\end{equation}
Its image in $M$ is $f^{-1}(J)$. A full kernel reduction tracks column
operations in an auxiliary matrix initialized to the identity
\cite[Section~3]{mpfree}. Since only the first $n$ coordinates are needed,
we instead initialize the auxiliary matrix to $\pi$ and apply the same
column operations to it. When a column of $D$ becomes zero, its auxiliary
column is returned if nonzero. Column permutations and grade assignments
are shared by the two matrices.

\begin{proposition}
The projected reduction returns graded generators of $\pi(\ker D)$.
\end{proposition}
\begin{proof}
Let $A$ denote the operation history of the corresponding full kernel
reduction, so its current matrix is $DA$. The projected auxiliary matrix
is initially $\pi$ and remains $\pi A$, since each column operation
preserves this identity. Thus every kernel generator $k$ emitted by the
full reduction is replaced by $\pi k$ at the same grade. If the full
generators are $k_1,\ldots,k_s$, their images generate $\pi(\ker D)$;
zero images may be discarded.
\end{proof}

This avoids storing and adding the witness coordinates $y,z$. For $q$
columns of $D$, the auxiliary matrix has $n$ rather than $q$ rows; the
reduction of $D$ itself is unchanged. The gain depends on the nonzero
entries in the discarded auxiliary rows. Projected generators can still
be redundant and may require subsequent minimization.

\subsection{Restricting I preimages and reusing their target}
Two changes simplify the repeated I updates. First, compute the preimage
on the current submodule rather than on all of $M$. If
$G:P\twoheadrightarrow U_{i-1}$ is a free generator map, then
\begin{equation}\label{eq:fused}
 U_{i-1}\cap f_i^{-1}(J)=G\bigl((f_iG)^{-1}(J)\bigr).
\end{equation}
Second, reduce the common target generators once per pass. If $S$ lifts
generators of $J$, replace $C=[\,R_N\mid S\,]$ by a graded generator
matrix $C'$ with the same image in the same free target module. Each
preimage then uses $C'$ as its witness block.

\begin{proposition}
Equation~\eqref{eq:fused} holds without assuming that $G$ is injective.
Moreover, for any graded map $A$ into the common free target,
\begin{equation}\label{eq:target}
 \pi_A\ker[\,A\mid C\,]=\pi_A\ker[\,A\mid C'\,],
\end{equation}
where $\pi_A$ projects onto the domain of $A$.
\end{proposition}
\begin{proof}
An element $x$ of the left side of~\eqref{eq:fused} has a lift $p\in P$
with $Gp=x$ and $f_iGp\in J$; conversely, every such $p$ maps into
that intersection. For~\eqref{eq:target}, in each grade a vector $x$
belongs to the first projected kernel exactly when $Ax\in\im C$.
Since $\im C=\im C'$ as graded submodules, this is equivalent to
membership in the second projected kernel.
\end{proof}

The fused update multiplies the lift of $f_i$ by $G$, computes the
preimage in $P$, and maps its generators back by $G$. It eliminates a
separate intersection and is particularly useful when the current
generator domain is small. Target preparation removes redundant witness
columns from every preimage in the pass, although each map still requires
its own kernel reduction. Both changes combine with projected tracking.
The target reduction must preserve the ambient row basis and the image
at every grade; presentation minimization that cancels rows cannot be
substituted without also transforming the maps.

\subsection{Deferring reductions in the K sum}
If $G$ generates the fixed submodule $V$, the columns of $h_jG$ already
generate $h_j(V)$. We can therefore omit their individual minimization
and append them directly to the accumulator for $L$. Independently, we
can minimize the accumulated sum after each batch of $b$ images instead
of after every image; the implementation uses $b=16$ and also reduces
the final partial batch.

\begin{proposition}
Omitting individual image minimizations and batching sum minimizations
preserve $L$ and $V_{\mathrm{next}}$, separately or together.
\end{proposition}
\begin{proof}
Appending generator matrices represents the sum of their images, and
exact generator minimization preserves that image. Hence, after processing
the first $k$ maps, the accumulator generates
$\eta V+\sum_{j=1}^{k}h_j(V)$, regardless of the reduction schedule.
As $V$ is fixed, the final sum and its preimage under $\eta$ are unchanged.
\end{proof}

For nonzero $V$ and $r$ images, these changes remove $r$ individual image reductions and
replace $r$ sum reductions by $\lceil r/b\rceil$. This excludes the
initial construction of $\eta V$ and the final preimage. Batching trades
fewer reductions for larger intermediate matrices. If the map family or
$V$ is empty, the preimage step still runs: in particular, at $V=0$ it
returns $\ker\eta$.

\subsection{Pivot lookup and progress output}
The reduction stores the owner of each pivot row in a row-indexed array
instead of a balanced search tree. Both structures represent the same
partial function from rows to columns, so, by induction over lookups and
updates, they choose the same reducers and perform the same column
operations. Array access takes constant time, at the cost of initialization
and storage proportional to the number of rows. Finally, reporting progress
only when the integer percentage changes bounds the number of writes per
pass by 101. These writes do not enter any algebraic computation.

\begin{thebibliography}{1}
\bibitem{mpfree}
M.~Kerber and A.~Rolle.
\emph{Fast Minimal Presentations of Bi-graded Persistence Modules}.
ALENEX 2021, pp.~207--220.
\href{https://arxiv.org/abs/2010.15623}{arXiv:2010.15623}.
\end{thebibliography}
\end{document}
